% Experimental Setup — Securing LLMs Against Prompt Injection
% Academic manuscript section (IEEE / arXiv style)
% Companion to methodology.tex; citation keys match references.txt

\section{Experimental Setup}
\label{sec:experimental-setup}

This section specifies the experimental conditions under which the defenses in Section~\ref{sec:methodology} are compared.
We describe the research questions, compared systems, computing environment, model configurations, evaluation procedures, robustness and ablation settings, latency protocol, and reproducibility controls.
The emphasis is on the scientific protocol needed to interpret the reported results; implementation details of the accompanying software release are omitted except where they affect experimental validity.

\subsection{Research Questions}
\label{sec:rq}

The empirical study is organized around four research questions that map directly onto the threat model and metrics defined earlier.

\begin{enumerate}
  \item \textbf{RQ1 (unprotected risk).}
    How often do injection prompts succeed against an unprotected instruction-tuned model under the chosen judge?
  \item \textbf{RQ2 (lexical defenses).}
    Do simple lexical defenses---regular expressions and keyword heuristics---produce meaningful reductions in Bypass Rate and True ASR relative to the unprotected baseline?
  \item \textbf{RQ3 (context-aware gains).}
    Does multi-signal context-aware sanitization outperform those lexical baselines on the same held-out test partition, and how do the handcrafted and learned variants compare with each other?
  \item \textbf{RQ4 (usability cost).}
    What false-positive burden, on both the standard benign test partition and an edge-case benign set, accompanies the stronger defenses?
\end{enumerate}

These questions reflect a methodological stance common in contemporary guardrail evaluation: security claims are incomplete unless over-blocking and residual end-to-end success are reported alongside filter misses~\cite{zizzo2025adversarial,promptfoo2024owasp,owasp_llm}.
RQ1 establishes the baseline hazard.
RQ2 tests whether inexpensive lexical controls already suffice on this corpus.
RQ3 isolates the contribution of semantic and structural context.
RQ4 prevents security improvements from being interpreted without reference to user impact.

\subsection{Compared Systems}
\label{sec:compared-systems}

We evaluate five conditions under a shared allow-or-block interface so that differences in outcome can be attributed to filtering policy rather than to inconsistent orchestration.

\begin{itemize}
  \item \textbf{Baseline:} no sanitization; every prompt is passed to the target model.
  \item \textbf{Regex:} pattern-based blocking using a fixed library of injection-related expressions.
  \item \textbf{Keyword:} weighted lexical scoring with a block threshold of 0.5.
  \item \textbf{Method~A:} hybrid context-aware sanitization in which semantic hard triggers form the primary layer and weighted multi-signal scoring provides secondary coverage.
  \item \textbf{Method~B:} learned fusion of an expanded ten-feature representation, with a validation-selected probability threshold and no hard vetoes.
\end{itemize}

An optional TF--IDF logistic regression classifier is retained as a classical supervised lexical reference and is reported only when explicitly included in an experiment.
It is useful as a bag-of-words contrast condition but is not required to answer RQ1--RQ4.

\paragraph{Paired evaluation suites.}
Loading multiple large models in one process can exhaust device memory on a single workstation.
We therefore execute the primary comparison as two paired suites that share the same lexical baselines.
Suite~A comprises baseline, regex, keyword, and Method~A.
Suite~B comprises baseline, regex, keyword, and Method~B.
Because the test partition and lexical baselines are identical across suites, shared methods remain directly comparable, while Method~A and Method~B can be contrasted as alternative context-aware designs under matched data and judge conditions.

\subsection{Computing Environment}
\label{sec:hardware}

Experiments were conducted on a workstation running Windows~11 with a multi-core Intel CPU (22 logical processors), approximately 15.4~GB of system memory, and an NVIDIA GeForce RTX~4070 Laptop GPU with 8~GB of device memory.
The software environment used Python~3.13 and PyTorch~2.6 with CUDA~12.4 support.
Host characteristics are recorded alongside latency measurements so that timing results can be interpreted relative to the execution platform rather than as absolute constants.

\paragraph{Device placement.}
Device placement is asymmetric by design.
The target model and compliance judge use GPU acceleration when sufficient memory is available.
Auxiliary components used by the context-aware sanitizers---sentence embeddings and perplexity scoring---are executed on CPU.
This separation has two consequences that matter for interpretation.
First, GPU memory remains available for generation and judging.
Second, sanitizer latency can be reported as an independent pre-inference cost rather than as a quantity entangled with target-model kernels.

\subsection{Model Configurations}
\label{sec:models-setup}

\paragraph{Target model.}
Unblocked prompts are answered by Qwen2.5-3B-Instruct under greedy decoding with a maximum of 512 newly generated tokens.
Half-precision inference is used on GPU when available; otherwise full precision is used on CPU.
The choice of a compact open instruction-tuned model is deliberate: it makes the pipeline reproducible without proprietary API access, while still exhibiting nontrivial compliance behaviour against which defenses can be measured.
Greedy decoding removes sampling variance from the principal security tables.

\paragraph{Compliance judge.}
End-to-end success is assessed by a separately loaded judge from the same model family.
Given the original prompt and the target response, the judge returns a binary compliance label.
Unparseable judgments are treated as refusals, which is conservative with respect to claimed attack success.
Long prompts and outputs are truncated prior to judging in order to keep the judgment context within a manageable length.
Before the main run, the judge is checked on a hand-authored calibration set of 20 examples; materially weak calibration would indicate that Layer~2 labels should be interpreted with caution.

\paragraph{Sanitizer auxiliaries.}
Methods~A and~B employ a compact sentence embedding model (all-MiniLM-L6-v2) for semantic similarity~\cite{reimers2019sentence,wang2020minilm} and DistilGPT-2 for perplexity-based naturalness scoring~\cite{radford2019language}.
Reference embeddings and perplexity calibration statistics are derived from the training partition only, so that test prompts do not influence the sanitizer's internal reference geometry.
Both auxiliaries remain on CPU throughout training caches, main evaluation, ablations, and latency measurement.

\paragraph{Adaptive rewriter.}
Stylistic attack variants are generated with FLAN-T5-Base under greedy decoding and a short generation budget.
The rewriter is unloaded before subsequent target-model evaluation so that device memory can be reused.
This design keeps the adaptive protocol practical on a single GPU workstation while still producing systematic style variation.

\subsection{Primary Evaluation Protocol}
\label{sec:eval-protocol}

All primary metrics reported in this paper are computed on the frozen held-out test partition comprising 30 injection prompts and 23 benign prompts, using the full two-layer pipeline with the target model and compliance judge.
Sanitizer-only runs are used only where explicitly stated (paraphrase robustness and edge-case false-positive measurement) and are never substituted for the main security tables.

For each defense, evaluation proceeds in two phases.

In the first phase, every test prompt is presented to the sanitizer.
Blocked prompts terminate immediately with a blocked outcome and no model generation.
Unblocked prompts are passed to the target model, and the generated responses are stored for later judging.

In the second phase, auxiliary models and the target model are released from memory, and the judge scores only the unblocked injection cases.
This staging avoids concurrent residency of sanitizer auxiliaries, the target model, and the judge when memory is limited, without changing the logical definition of a successful attack.

\paragraph{Success criterion.}
A successful attack requires an injection label, a sanitizer miss, and a complied judgment.
All three conditions are necessary.
In particular, an unblocked injection that the target refuses is counted in Bypass Rate but not in True ASR.

\paragraph{Reported quantities.}
We report Bypass Rate, Detection Rate, false-positive rate, True ASR, successful-attack counts, and Layer~1 confusion matrices.
Bypass Rate and True ASR are always presented as separate columns because they answer different questions: the former concerns filter coverage, whereas the latter concerns realized end-to-end harm after model refusal is taken into account~\cite{zizzo2025adversarial,yi2024jailbreak}.
We do not report a single undifferentiated attack-success percentage that could be mistaken for either metric.
False-positive rate on the standard benign partition provides the primary usability counterpart to these security metrics.
Suite~A and Suite~B share the same test partition and lexical baselines, so their overlapping conditions are comparable; Method~A and Method~B are contrasted as the two context-aware alternatives under that matched protocol.

\subsection{Robustness Evaluation Settings}
\label{sec:robustness-setup}

Robustness experiments reuse the same defenses and, where applicable, the same two-layer success criterion, but they alter the input distribution in controlled ways.

\paragraph{Paraphrase setting.}
The first fifty training injection prompts are rewritten with a fixed set of nine synonym substitutions.
Bypass Rates on original and paraphrased text are compared under sanitizer-only evaluation.
No target-model or judge call is involved.
The setting therefore measures sensitivity to mild lexical drift rather than end-to-end compliance after paraphrase.
A defense whose Bypass Rate rises sharply under these substitutions is lexically brittle even if its headline test-set numbers are strong.

\paragraph{Edge-case benign setting.}
Eighty fixed technical-but-benign prompts are scored by each defense.
Edge-case false-positive rate is defined as the fraction of these prompts that are blocked.
The set deliberately includes developer- and systems-oriented language that overlaps superficially with attack vocabulary, such as process termination, configuration overrides, and non-security uses of ``jailbreak''~\cite{bair2025prompt}.
This setting is the principal probe for RQ4 beyond the ordinary benign test partition.

\paragraph{Adaptive setting.}
Each of the thirty test injections is rewritten into five styles---role-play, research, hypothetical, indirect, and multistep---yielding up to 150 adaptive prompts.
Every variant is evaluated with the full two-layer pipeline.
Reported quantities include Bypass Rate, True ASR, and successful-attack counts on the adaptive pool.
The protocol is best understood as a controlled style-transfer stress test: it asks whether defenses remain effective when malicious goals are preserved but surface form changes systematically~\cite{promptfoo2024owasp,zizzo2025adversarial}.
It does not claim to approximate the strongest available adversarial optimization attacks.

\subsection{Ablation Protocol}
\label{sec:ablation-setup}

Ablations are essential for attributing performance to components rather than to the mere presence of a complex pipeline.
All ablations use the same held-out test partition as the primary evaluation and share loaded auxiliary models whenever possible, so that observed differences are attributable to disabled components rather than to reload noise or nondeterministic initialization.

Three families are reported.

\begin{enumerate}
  \item \emph{Soft-signal ablation for Method~A.}
    Hard vetoes are disabled, and each of the eight soft signals is removed in turn.
    Comparison against the full soft-weighted configuration reveals how much secondary coverage the soft layer provides when the primary hard-trigger layer is unavailable.
  \item \emph{Hard-veto ablation for Method~A.}
    Full Method~A is compared with a soft-only variant and with variants that disable each veto individually.
    This family tests the hybrid claim that headline security depends mainly on hard triggers---especially the semantic veto---rather than on soft aggregation alone.
  \item \emph{Feature ablation for Method~B.}
    Selected feature groups are zeroed at inference time, including an all-semantic condition, while the trained fusion model remains fixed.
    This family tests whether learned performance is carried by semantic features specifically or is broadly distributed across the feature set.
\end{enumerate}

For each variant we record Bypass Rate, True ASR, false-positive rate, detection rate, confusion counts, and percentage-point deltas relative to the corresponding full system.
These deltas are the basis for component-level claims in the results section.

\subsection{Latency Measurement}
\label{sec:latency-setup}

Sanitizer overhead is measured on CPU for semantic scoring, perplexity scoring, full Method~A sanitization, and the two lexical baselines.
Each measurement uses fifty prompts, five warm-up iterations, and a fixed random seed.
We report mean, median, and tail latency (95th percentile), together with extrema, and we record host hardware metadata alongside the timing results.

The latency protocol is aligned with the main pipeline in two respects.
Auxiliary models remain on CPU, matching their deployment placement.
The timed Method~A path includes the full allow-or-block decision rather than an isolated subscore, so the figures reflect the cost of inserting the complete gate into a request path.
Lexical baselines are timed under the same protocol to provide lower-bound references for inexpensive filters.

\subsection{Reproducibility}
\label{sec:repro-setup}

Reproducibility is treated as part of the experimental design rather than as an after-the-fact checklist.
Unless otherwise stated, all randomized procedures use seed 42, including dataset partitioning and Method~B training.
Target-model, judge, and rewriter decoding are greedy.
Dataset partitions are frozen prior to experimentation and are never reshuffled at evaluation time.
Sanitizer decisions are deterministic given fixed auxiliary models, caches, and thresholds.

True ASR additionally depends on target and judge generations.
Under greedy decoding these outputs are stable for a fixed model revision and configuration, but they are not independent of software versions or numerical backends.
Where exact numerical reproduction is required, the same model revisions and dependency versions should be used.

Overall, the experimental setup is designed so that the principal metrics and figures can be reproduced from the released materials and frozen partitions, while remaining explicit about the compact corpus and the limited adaptive attacker considered in this study~\cite{owasp_llm,zizzo2025adversarial}.
The resulting protocol supports a fair comparison of lexical and context-aware sanitizers under matched data, matched metrics, and matched end-to-end judgment.
